\chapter{Continuous solutions and normalization}
\label{chap:weak-equations}

The bounded product from \cref{chap:products} makes it possible to formulate equations with nonsmooth densities on an arbitrary compact Hermitian manifold. We prove continuous existence from smooth solvability, local capacity estimates and modified comparison. The only imported analytic estimates in this chapter are classical local Bedford--Taylor estimates. We also distinguish uniqueness of a scaling constant from uniqueness of a potential.

\section{The two basic existence theorems}

Let $X$ be a connected compact complex manifold of dimension $n\geq1$, let $\omega$ be a Hermitian form, and let $\mathrm{d}V$ be a smooth positive volume form on $X$. For $u\in\PSH(X,\omega)$ write $\omega_u=\omega+\ddc u$; products of bounded potentials have the meaning of \cref{chap:products}. Fix $p>1$ and a nonnegative $f\in L^p(X,\mathrm{d}V)$ with $\int_Xf\,\mathrm{d}V>0$. Multiplying by the smooth positive ratio $\mathrm{d}V/\omega^n$ converts this setting into the one used in the references.

\begin{theorem}[Ko{\l}odziej--Nguyen]\label{thm:kn-continuous}
There are $c>0$ and $u\in\PSH(X,\omega)\cap C(X)$ such that
\[
 (\omega+\ddc u)^n=c f\,\mathrm{d}V.
\]
One may normalize $\sup_Xu=0$.
\end{theorem}

This is \cite[Theorem 0.1]{KN15}. It requires neither the bounded mass property nor a positive lower-volume hypothesis. Uniqueness of the normalized potential fails in this generality, even for a smooth density and a background with bounded mass; see \cref{thm:degenerate-hermitian-nonuniqueness}.

\begin{theorem}[Nguyen]\label{thm:nguyen-exp}
For every $\lambda>0$, there is a unique continuous $\omega$-plurisubharmonic function $u_\lambda$ satisfying
\begin{equation}\label{eq:weak-exp}
 (\omega+\ddc u_\lambda)^n=\mathrm{e}^{\lambda u_\lambda}f\,\mathrm{d}V.
\end{equation}
\end{theorem}

This is \cite[Theorem 0.1]{Nguyen16}. In both references $\ddc=\mathrm i\partial\bar\partial/\pi$, twice our operator. Replacing a reference potential $\varphi$ by $u=2\varphi$ gives the stated current. For \eqref{eq:weak-exp}, use the reference parameter $2\lambda$ as well. Thus the change of normalization also affects the exponential parameter.

The assumption $\lambda>0$ removes the additive freedom: shifting $u_\lambda$ changes the right side. These theorems concern strictly positive smooth backgrounds and continuous potentials. They should not be applied to a merely big background or an arbitrary singular potential without a further theorem.

\section{The modified comparison principle}

The variation of total mass prevents the ordinary global comparison principle. Replacing a comparison potential $v$ by $(1-\varepsilon)v$ reserves strict positivity that absorbs the torsion terms. On a sublevel of height $s$, integration by parts makes the error proportional to $s$; a finite induction controls the intervening mixed products. The multiplicative correction tends to one as the height tends to zero with $\varepsilon$ fixed.

Choose $B_0\geq1$ such that
\[
 -B_0\omega^2\leq\ddc\omega\leq B_0\omega^2,
 \qquad
 -B_0\omega^3\leq \mathrm{d}\omega\wedge\dc\omega\leq B_0\omega^3.
\]
Only inequalities after wedging with positive complementary forms are needed. Such a constant exists because $\omega$ is Hermitian and $X$ is compact. Terms whose degree exceeds the dimension are omitted.

\begin{theorem}\label{thm:modified-comparison}\index{Comparison principle!modified}
Let $u,v\in\PSH(X,\omega)\cap L^\infty(X)$ and $0<\varepsilon<1$. For $s>0$, put
\[
 m_\varepsilon=\inf_X\bigl(u-(1-\varepsilon)v\bigr),\qquad
 U_{\varepsilon,s}=\{u<(1-\varepsilon)v+m_\varepsilon+s\}.
\]
If $0<s<\varepsilon^3/(32B_0n^2)$, then, with
$a=B_0n^2s\varepsilon^{-3}$,
\begin{equation}\label{eq:modified-comparison}
 (1-4a)^n\int_{U_{\varepsilon,s}}\omega_{(1-\varepsilon)v}^{\,n}
 \leq\int_{U_{\varepsilon,s}}\omega_u^n.
\end{equation}
In particular, the reverse factor $(1-4a)^{-n}$ tends to one as $s\downarrow0$ for fixed $\varepsilon$.
\end{theorem}
\begin{proof}
Write $v_\varepsilon=(1-\varepsilon)v$, $U=U_{\varepsilon,s}$ and
\[
 \phi=\max(u,v_\varepsilon+m_\varepsilon+s),\qquad
 T_k=\omega_u^k\wedge\omega_\phi^{n-k},\qquad I_k=\int_U T_k\quad(0\leq k\leq n).
\]
Set $I_{-1}=I_{-2}=0$. Then $0\leq\phi-u\leq s$, this difference vanishes outside $U$, and on $U$ locality gives
$\omega_\phi=\omega_{v_\varepsilon}\geq\varepsilon\omega$ in all bounded products.

By \cref{lem:bounded-ordered-contact}, for every fixed positive mixed product $W$ of $n-1$ bounded-potential factors,
\[
 \mathbf1_{\{u=\phi\}}\omega_u\wedge W
 \leq\mathbf1_{\{u=\phi\}}\omega_\phi\wedge W.
\]

For $0\leq k\leq n-1$, apply it with $W=\omega_u^k\wedge\omega_\phi^{n-k-1}$. On $X\setminus U$ the difference $T_k-T_{k+1}$ is nonnegative. Consequently
\begin{equation}\label{eq:comparison-stokes}
 I_k-I_{k+1}
 \leq\int_X\ddc(\phi-u)\wedge W
 =\int_X(\phi-u)\ddc W.
\end{equation}
This Stokes identity is valid for the bounded potentials in use. To check it, approximate $u$ and $\phi$ by bounded smooth decreasing $\omega$-plurisubharmonic functions using \cref{thm:bounded-exact-regularization}. For smooth functions it is ordinary integration by parts. The derivative $\ddc W$ expands as below in bounded Hessian products and smooth forms; bounded Bedford--Taylor convergence, also against bounded quasi-continuous coefficients converging in capacity, passes each term to the limit.

\Needspace{11\baselineskip}
Set $r=n-k-1$. Since $\mathrm{d}\omega_u=\mathrm{d}\omega_\phi=\mathrm{d}\omega$, expansion gives
\[
 \begin{split}
 \ddc W={}&k\ddc\omega\wedge\omega_u^{k-1}\wedge\omega_\phi^r
       +r\ddc\omega\wedge\omega_u^k\wedge\omega_\phi^{r-1}\\
 &+k(k-1)\mathrm{d}\omega\wedge\dc\omega\wedge\omega_u^{k-2}\wedge\omega_\phi^r\\
 &+2kr\,\mathrm{d}\omega\wedge\dc\omega\wedge\omega_u^{k-1}\wedge\omega_\phi^{r-1}\\
 &+r(r-1)\mathrm{d}\omega\wedge\dc\omega\wedge\omega_u^k\wedge\omega_\phi^{r-2}.
 \end{split}
\]
Factors with a zero coefficient and a negative power are absent. On $U$ we can replace each $\omega^j$ in an upper bound by $\varepsilon^{-j}\omega_\phi^j$. Using $\phi-u\leq s$ and $\varepsilon^{-2}\leq\varepsilon^{-3}$ in \eqref{eq:comparison-stokes}, we obtain
\[
 I_k-I_{k+1}\leq B_0s\varepsilon^{-3}
 \bigl(k(k-1)I_{k-2}+k(1+2r)I_{k-1}+r^2I_k\bigr).
\]
Each coefficient is at most $n^2$. Hence
\begin{equation}\label{eq:comparison-recurrence}
 I_k-I_{k+1}\leq a(I_{k-2}+I_{k-1}+I_k).
\end{equation}
Put $q=1-4a>7/8$. Inductively suppose $qI_j\leq I_{j+1}$ for $j<k$. Then
$I_{k-1}\leq q^{-1}I_k$ and $I_{k-2}\leq q^{-2}I_k$ whenever the indices exist. Since
\[
 1+q^{-1}+q^{-2}\leq1+8/7+64/49<4,
\]
\eqref{eq:comparison-recurrence} gives $qI_k\leq I_{k+1}$. The case $k=0$ follows from the same formula with the two absent terms set to zero. Induction proves $q^n I_0\leq I_n$. On $U$, $T_0=\omega_{v_\varepsilon}^n$, while $T_n=\omega_u^n$, giving the theorem. In dimension one all torsion terms vanish and the same argument applies.
\end{proof}

A strict multiplicative inequality between the measures absorbs the comparison error. Choose the positivity reserve first, and then a sublevel thin enough to retain this strict inequality.

\begin{theorem}\label{thm:bounded-strict-domination}
Let $u,v\in\PSH(X,\omega)\cap L^\infty(X)$, and let $0\leq c<1$. If
\[
 \mathbf1_{\{u<v\}}\omega_u^n
 \leq c\,\mathbf1_{\{u<v\}}\omega_v^n,
\]
then $u\geq v$. In particular, $\omega_u^n(\{u<v\})=0$ implies $u\geq v$.
\end{theorem}
\begin{proof}
Suppose $m=\inf_X(u-v)<0$. Choose $\varepsilon>0$ so small that
\[
 (1-\varepsilon)^n>c,
 \qquad 2\varepsilon\|v\|_\infty<-m/4,
\]
omitting the second restriction if $v=0$. Use $m_\varepsilon$ and $U_{\varepsilon,s}$ defined in \cref{thm:modified-comparison} for this pair $u,v$. Since
$|m_\varepsilon-m|\leq\varepsilon\|v\|_\infty$, for all sufficiently small $s>0$ the set $U=U_{\varepsilon,s}$ is contained in $\{u<v\}$: there
\[
 u-v<m_\varepsilon+s-\varepsilon v
       \leq m+2\varepsilon\|v\|_\infty+s<0.
\]
The set $U$ is nonempty by the definition of the infimum and has positive volume by \cref{lem:strict-set-positive-volume}.

\Needspace{8\baselineskip}
Make $s$ smaller so that \cref{thm:modified-comparison} applies and
$(1-4B_0n^2s\varepsilon^{-3})^n(1-\varepsilon)^n>c$. Write the first factor as $q^n$. Positivity of every mixed term gives
\[
 \omega_{(1-\varepsilon)v}^n
 \geq(1-\varepsilon)^n\omega_v^n+\varepsilon^n\omega^n.
\]
Comparison and the assumption would therefore imply
\[
 q^n(1-\varepsilon)^n\int_U\omega_v^n
       +q^n\varepsilon^n\int_U\omega^n
 \leq\int_U\omega_u^n\leq c\int_U\omega_v^n.
\]
All integrals are finite. The chosen strict inequality and the positive volume of $U$ make this impossible.
\end{proof}

\begin{corollary}\label{cor:bounded-mass-positive}
Every bounded $\omega$-plurisubharmonic function has strictly positive Monge--Amp\`ere mass.
\end{corollary}
\begin{proof}
Let $u$ be a bounded $\omega$-plurisubharmonic function. If $\omega_u^n=0$, bounded domination with the constant potential $v=\sup_Xu+1$ would give $u\geq v$, a contradiction.
\end{proof}

The lower mass bound here depends on the potential.

\Needspace{8\baselineskip}
\begin{corollary}\label{cor:bounded-exponential-comparison}\index{Comparison principle!exponential}
If $u,v\in\PSH(X,\omega)\cap L^\infty(X)$ and $\lambda>0$, then
\[
 \mathrm{e}^{-\lambda v}\omega_v^n\geq \mathrm{e}^{-\lambda u}\omega_u^n
 \quad\Longrightarrow\quad v\leq u.
\]
\end{corollary}
\begin{proof}
For every $t>0$, on $\{u<v-t\}$ the assumption gives
$\omega_u^n\leq \mathrm{e}^{-\lambda t}\omega_v^n$. Apply \cref{thm:bounded-strict-domination} to $u,v-t$, and then let $t\downarrow0$.
\end{proof}


\section{Capacity estimates and uniform compactness}

Insert a capacity test into the comparison potential. Modified comparison then bounds the capacity of one sublevel by the Monge--Amp\`ere mass of a slightly larger one. For an $L^p$ density, the capacity estimate of \cref{cor:lp-capacity-square} reverses this relation: a nonempty sublevel of fixed height has a definite capacity. The tail bound \eqref{eq:global-capacity-tail} bounds the depth of a normalized potential. Applied to differences of solutions, the same estimate upgrades $L^1$ compactness to uniform compactness. Write $\operatorname{Cap}=\operatorname{Cap}_\omega$ for the capacity of \cref{sec:hermitian-capacity}.

\Needspace{15\baselineskip}
\begin{lemma}\label{lem:thin-sublevel-capacity}
Let $u,v\in\PSH(X,\omega)\cap L^\infty(X)$, with $-H\leq v\leq0$ and $H\geq1$. Fix $0<\varepsilon<1$, and put
\[
 m=\inf_X(u-(1-\varepsilon)v),\qquad
 U_s=\{u<(1-\varepsilon)v+m+s\}.
\]
There are $\delta_\varepsilon>0$, depending only on $(X,\omega,\varepsilon)$, and $C_H>0$, depending only on $(n,H)$, such that
\begin{equation}\label{eq:thin-capacity}
 t^n\operatorname{Cap}(U_s)\leq C_H\int_{U_{s+t}}\omega_u^n
 \qquad(0<s,t,\ s+t<\delta_\varepsilon).
\end{equation}
Every $U_s$, $s>0$, has positive capacity.
\end{lemma}
\begin{proof}
Use the constant $B_0$ from \cref{thm:modified-comparison} and set
\[
 \delta_\varepsilon=\min\{1-\varepsilon,\varepsilon^3/(128B_0n^2)\}.
\]
For a capacity test $0\leq\rho\leq1$, put
\[
 \tau=\frac{t}{2(1-\varepsilon)(H+1)},\quad
 w=(1-\tau)v+\tau\rho,\quad
 m'=\inf_X(u-(1-\varepsilon)w).
\]
Then $0<\tau<1$ and
$0\leq(1-\varepsilon)(w-v)\leq t/2$, so $m-t/2\leq m'\leq m$. Consequently
\[
 U_s\subset V:=\{u<(1-\varepsilon)w+m'+s+t/2\}
 \subset U_{s+t}.
\]
Apply modified comparison to $u,w$ at height $s+t/2$. Its inverse factor is at most $(8/7)^n$. Moreover,
\[
 \omega_{(1-\varepsilon)w}
 \geq(1-\varepsilon)\tau\omega_\rho
 =\frac{t}{2(H+1)}\omega_\rho.
\]
Integration over $U_s\subset V$ and comparison over $V$ give \eqref{eq:thin-capacity}, with $C_H=[2(H+1)]^n(8/7)^n$, after taking the supremum over $\rho$. Each $U_s$ is nonempty by the definition of the infimum. \Cref{lem:strict-set-positive-volume} gives positive smooth volume, hence positive capacity by the constant test.
\end{proof}

\Needspace{11\baselineskip}
\begin{lemma}\label{lem:backwards-capacity}
In the preceding setting, suppose
\[
 \int_E\omega_u^n\leq A\operatorname{Cap}(E)^2\qquad(E\subset X\ \hbox{Borel}),
 \qquad A>0.
\]
For every $0<D<\delta_\varepsilon$,
\begin{equation}\label{eq:sublevel-capacity-lower}
 \operatorname{Cap}(U_D)\geq\frac{D^n}{4^nC_HA}.
\end{equation}
\end{lemma}
\begin{proof}
Write $a(s)=\operatorname{Cap}(U_s)^{1/n}>0$ and $K=(C_HA)^{1/n}$. The preceding lemma gives
\[
 t\,a(s)\leq K a(s+t)^2\qquad(s+t<\delta_\varepsilon).
\]
Suppose $4Ka(D)<D$. Starting with $s_0=D$, define
\[
 t_j=2Ka(s_j),\qquad s_{j+1}=s_j-t_j.
\]
We check positivity before using the sublevel inequality. If $a(s_i)\leq2^{-i}a(D)$ for $0\leq i\leq j$, then
\[
 \sum_{i=0}^j t_i\leq2Ka(D)\sum_{i=0}^j2^{-i}<4Ka(D)<D.
\]
Hence $s_{j+1}>0$, and the displayed sublevel inequality gives $a(s_{j+1})\leq a(s_j)/2$. Induction justifies every step, and $s_j\downarrow s_\infty\geq D-4Ka(D)>0$. Monotonicity gives
$0<a(s_\infty)\leq a(s_j)\leq2^{-j}a(D)$, a contradiction. Therefore $D\leq4Ka(D)$, which is \eqref{eq:sublevel-capacity-lower}.
\end{proof}

\begin{theorem}\label{thm:lp-potential-estimate}
There are constants $C_p,c_p>0$, depending only on $(X,\omega,\mathrm{d}V,p)$, such that, whenever
$u\in\PSH(X,\omega)\cap L^\infty(X)$, $\sup_Xu=0$, and
\[
 \omega_u^n=F\,\mathrm{d}V,\qquad F\in L^p(X,\mathrm{d}V),\quad F\geq0,
\]
we have
\begin{equation}\label{eq:lp-potential-bound}
 \|u\|_\infty\leq C_p(1+\|F\|_p),
 \qquad \|F\|_p\geq c_p.
\end{equation}
\end{theorem}
\begin{proof}
The zero density is excluded by \cref{cor:bounded-mass-positive}. Apply \cref{cor:lp-capacity-square,lem:backwards-capacity} with $v=0$, $H=1$, $\varepsilon=1/2$, and $A=C_p\|F\|_p>0$. Fix $D=\delta_{1/2}/2$. If $M=-\inf_Xu$, the resulting inequality is
\[
 \operatorname{Cap}(\{u<-M+D\})\geq c/\|F\|_p.
\]
If $M>D+1$, \eqref{eq:global-capacity-tail} bounds this capacity by $C/(M-D)$, giving the first assertion; the remaining case is immediate. Bounding the same capacity by $\operatorname{Cap}(X)$ gives the second assertion.
\end{proof}

If two uniformly bounded solutions remain separated at some point, the preceding estimate gives a fixed positive capacity to a sublevel of their difference. Small $L^1$ distance makes its volume small, while the uniform $L^p$ density bound makes its Monge--Amp\`ere mass small. This contradicts the sublevel estimate and upgrades $L^1$ compactness to uniform compactness. Put $p'=p/(p-1)$.

\begin{theorem}\label{thm:lp-uniform-compactness}
Suppose $u_j\in\PSH(X,\omega)\cap C(X)$, $\sup_Xu_j=0$, and
\[
 \omega_{u_j}^n=F_j\,\mathrm{d}V,\qquad F_j\geq0,\qquad \sup_j\|F_j\|_p<\infty
\]
for the fixed $p>1$. If $(u_j)$ is Cauchy in $L^1(X,\mathrm{d}V)$, then it is Cauchy uniformly. In particular, every such family has a subsequence converging uniformly to a continuous $\omega$-plurisubharmonic function.
\end{theorem}
\begin{proof}
By \cref{thm:lp-potential-estimate}, there is a common $H\geq1$ with $-H\leq u_j\leq0$. Suppose uniform Cauchy convergence fails. After exchanging indices when necessary, there are a number $a>0$ and arbitrarily large $j,k$ with
\[
 S:=\inf_X(u_k-u_j)\leq-a.
\]
\Needspace{10\baselineskip}
Choose $0<\varepsilon<1$ such that $\varepsilon H<a/4$, and choose a fixed $D>0$ with $2D<\delta_\varepsilon$ and $2D<a/4$. Apply the preceding sublevel lemmas with $u=u_k$, $v=u_j$. Since $v\leq0$,
$m=\inf(u-(1-\varepsilon)v)\leq S$. Thus
\[
 U_{2D}\subset\{u-v<S+2D+\varepsilon H\}
 \subset\{|u-v|>a/2\}.
\]
The capacity lower bound \eqref{eq:sublevel-capacity-lower}, with a common $A$ supplied by \cref{cor:lp-capacity-square}, gives $\operatorname{Cap}(U_D)\geq c>0$. On the other hand, \eqref{eq:thin-capacity}, H\"older and Chebyshev give
\[
 \begin{split}
 D^n\operatorname{Cap}(U_D)
 &\leq C_H\int_{U_{2D}}F_k\,\mathrm{d}V\\
 &\leq C_H\|F_k\|_p
       \left(\frac{2\|u_k-u_j\|_1}{a}\right)^{1/p'}
 \longrightarrow0.
 \end{split}
\]
This contradiction proves uniform Cauchy convergence. The normalized compactness theorem, \cref{thm:normalized-compactness}, supplies an $L^1$-Cauchy subsequence of any such family. Its uniform limit is continuous, and distributional convergence preserves $\omega$-plurisubharmonicity.
\end{proof}

\section{Construction of continuous solutions}

Smooth approximation gives solutions with unknown normalization scalars. A Gauduchon trace bounds those scalars above, and the lower bound for the density norm bounds them away from zero. The resulting uniform $L^p$ bound puts the potentials in the compact family just constructed, so their limit solves the prescribed equation.

\begin{proof}[of \cref{thm:kn-continuous}]
Choose smooth strictly positive $f_j\to f$ in $L^p(X,\mathrm{d}V)$. By \cref{thm:tw-smooth}, smooth existence gives
\[
 \omega_{u_j}^n=c_jf_j\,\mathrm{d}V,\qquad \sup_Xu_j=0,\qquad c_j>0.
\]
We first control the constants. For $n\geq2$, fix a Gauduchon metric $\gamma$. By \cref{prop:gauduchon-amgm},
\[
 c_j^{1/n}A_j\leq\int_X\omega\wedge\gamma^{n-1},
 \qquad A_j=\int_X(f_j\mathrm{d}V/\gamma^n)^{1/n}\gamma^n.
\]
Since $|x^{1/n}-y^{1/n}|\leq|x-y|^{1/n}$ for $x,y\geq0$, H\"older's inequality gives $A_j\to A>0$, where $A=\int_X(f\mathrm{d}V/\gamma^n)^{1/n}\gamma^n$. Thus $c_j$ is bounded above. In dimension one the same estimate is Stokes' identity and $A_j=\int_Xf_j\,\mathrm{d}V$.

The lower-density estimate in \cref{thm:lp-potential-estimate} gives
$c_j\|f_j\|_p\geq c_p$, so the constants are bounded away from zero. The densities $c_jf_j$ are now uniformly bounded in $L^p$, and the same theorem gives a common bound for $u_j$. Pass to a subsequence with $c_j\to c>0$ and $u_j$ Cauchy in $L^1$. By \cref{thm:lp-uniform-compactness}, $u_j\to u$ uniformly, with $u$ continuous and $\sup_Xu=0$.

Bounded Bedford--Taylor continuity gives $\omega_{u_j}^n\rightharpoonup\omega_u^n$, while $c_jf_j\to cf$ in $L^p$. Hence $\omega_u^n=cf\,\mathrm{d}V$, as required.
\end{proof}

For the exponential equation, the corresponding scalar is $\mathrm{e}^{\lambda\sup_Xu}$. Its upper bound comes from Jensen's inequality: weighting a Gauduchon probability by the root of the density relates the conserved trace to $\sup_Xu$. The lower-density estimate again gives the reverse bound.

\begin{lemma}\label{lem:exponential-gauduchon-mean}
Fix $\lambda>0$. Suppose smooth $u_j$ with $\omega_{u_j}>0$ solve
\[
 \omega_{u_j}^n=\mathrm{e}^{\lambda u_j}f_j\,\mathrm{d}V,
\]
where smooth strictly positive $f_j$ converge in $L^p$ to a nonnegative nonzero $f$. Then $M_j=\sup_Xu_j$ is bounded above and below, uniformly in $j$.
\end{lemma}
\begin{proof}
Absorb the fixed smooth positive ratio $\mathrm{d}V/\omega^n$ into both $f_j$ and $f$. This preserves their $L^p$ convergence, so we may use $\mathrm{d}V=\omega^n$ in the calculation. Choose $\rho>0$ smooth with $\ddc(\rho\omega^{n-1})=0$, using \cref{thm:gauduchon}; for $n=1$, take $\rho=1$. Put
\[
 \alpha=\int_X\rho\omega^n,\quad
 A_j=\int_X f_j^{1/n}\rho\omega^n,\quad
 \psi_j=u_j-M_j\leq0.
\]
Set $A=\int_X f^{1/n}\rho\omega^n$. The root inequality used in the preceding existence proof gives $A_j\to A>0$. Arithmetic--geometric mean and the conserved trace imply
\[
 \alpha=\int_X\omega_{u_j}\wedge\rho\omega^{n-1}
 \geq \mathrm{e}^{\lambda M_j/n}\int_X \mathrm{e}^{\lambda\psi_j/n}f_j^{1/n}\rho\omega^n.
\]
Apply Jensen's inequality to the probability measure
$A_j^{-1}f_j^{1/n}\rho\omega^n$. Taking logarithms yields
\begin{equation}\label{eq:gauduchon-exponential-mean}
 \lambda M_j\leq n\log(\alpha/A_j)
 +\frac{\lambda}{A_j}\int_X(-\psi_j)f_j^{1/n}\rho\omega^n.
\end{equation}
The normalized potentials $\psi_j$ have uniform $L^q$ bounds for every finite $q$, by \cref{lem:uniform-exponential}. With $q=np/(np-1)$, H\"older gives a uniform bound for the last integral, since
$\|f_j^{1/n}\|_{np}=\|f_j\|_p^{1/n}$. This proves the upper bound for $M_j$.

Finally, the density $F_j=\mathrm{e}^{\lambda u_j}f_j$ satisfies
$\|F_j\|_p\leq \mathrm{e}^{\lambda M_j}\|f_j\|_p$. The lower-density bound of \cref{thm:lp-potential-estimate}, applied to $\psi_j$, gives $c_p\leq \mathrm{e}^{\lambda M_j}\|f_j\|_p$, which bounds $M_j$ below.
\end{proof}

\Needspace{9\baselineskip}
\begin{proof}[of \cref{thm:nguyen-exp}]
Choose smooth strictly positive $f_j\to f$ in $L^p$. The smooth proper equation, \cref{thm:smooth-proper-existence}, gives smooth $u_j$ satisfying the required equations. The preceding lemma bounds $M_j$ and consequently bounds $F_j=\mathrm{e}^{\lambda u_j}f_j$ in $L^p$. The normalized potentials $\psi_j=u_j-M_j$ are uniformly bounded by \cref{thm:lp-potential-estimate}. Pass to a subsequence with $M_j\to M$ and $\psi_j$ Cauchy in $L^1$. \Cref{thm:lp-uniform-compactness} gives uniform convergence $u_j\to u\in\PSH(X,\omega)\cap C(X)$.

Uniform convergence and $f_j\to f$ in $L^p$ give
$\mathrm{e}^{\lambda u_j}f_j\to \mathrm{e}^{\lambda u}f$ in $L^p$. Bounded-product continuity therefore gives $\omega_u^n=\mathrm{e}^{\lambda u}f\,\mathrm{d}V$. Bounded exponential comparison, \cref{cor:bounded-exponential-comparison}, gives uniqueness, in fact among all bounded solutions.
\end{proof}


\section{Order for the exponential equation}

Let $X$ be a connected compact complex manifold of dimension $n\geq1$, let $\omega$ be a Hermitian form, and let $\mathrm{d}V$ be a smooth positive volume form on $X$. For $u\in\PSH(X,\omega)$ write $\omega_u=\omega+\ddc u$; products of bounded potentials have the meaning of \cref{chap:products}. All $L^p$ spaces and almost-everywhere comparisons below are with respect to $\mathrm{d}V$.

\begin{lemma}\label{lem:exp-order}
Let $a,b\in\PSH(X,\omega)\cap L^\infty(X)$ solve
\[
 \omega_a^n=\mathrm{e}^{\lambda a}f_a\,\mathrm{d}V,\qquad
 \omega_b^n=\mathrm{e}^{\lambda b}f_b\,\mathrm{d}V,
\]
where $\lambda>0$ and $f_a,f_b$ are nonnegative nonzero $L^p$ densities, $p>1$. If $f_a\leq f_b$ almost everywhere, then $b\leq a$ on $X$.
\end{lemma}

\begin{proof}
The equations and $f_a\leq f_b$ give
$\mathrm{e}^{-\lambda b}\omega_b^n\geq \mathrm{e}^{-\lambda a}\omega_a^n$.
Apply \cref{cor:bounded-exponential-comparison} with $u=a$ and $v=b$.
\end{proof}

This recovers \cite[Lemma 2.3]{Nguyen16}. The proof only requires bounded potentials and the stated measure inequality; it does not require a positive lower bound for either density.

The order reverses because increasing the density can be offset by decreasing the potential in the exponential factor.

\section{Recovering the scaling constant}

Let $X$ be a connected compact complex manifold of dimension $n\geq1$, let $\omega$ be a Hermitian form, and let $\mathrm{d}V$ be a smooth positive volume form on $X$. For $u\in\PSH(X,\omega)$ write $\omega_u=\omega+\ddc u$; products of bounded potentials have the meaning of \cref{chap:products}. Fix $p>1$ and a nonnegative $f\in L^p(X,\mathrm{d}V)$ with $\int_Xf\,\mathrm{d}V>0$. For each $\lambda>0$, let $u_\lambda$ be the unique continuous solution of $\omega_{u_\lambda}^n=\mathrm{e}^{\lambda u_\lambda}f\,\mathrm{d}V$ supplied by \cref{thm:nguyen-exp}.

The exponential equation determines the constant in \cref{thm:kn-continuous} through a small-parameter limit. The following argument also gives a uniform error estimate.

\begin{proposition}\label{prop:vanishing-exponent}
Suppose $u\in\PSH(X,\omega)\cap C(X)$ and $c>0$ satisfy $\omega_u^n=cf\,\mathrm{d}V$. Set $m=\inf_Xu$ and $M=\sup_Xu$. Then the solution of \eqref{eq:weak-exp} obeys
\begin{equation}\label{eq:exp-barriers}
 u+\frac{\log c}{\lambda}-M
 \leq u_\lambda\leq
 u+\frac{\log c}{\lambda}-m.
\end{equation}
In particular,
\[
 \|\lambda u_\lambda-\log c\|_\infty
 \leq\lambda(M-m),
 \qquad
 \mathrm{e}^{\lambda u_\lambda}\longrightarrow c
 \quad\hbox{uniformly on }X.
\]
The scaling constant $c$ is unique.
\end{proposition}
\begin{proof}
Put $a=u+\lambda^{-1}\log c-M$ and $b=u+\lambda^{-1}\log c-m$. Both have the same Monge--Amp\`ere measure as $u$, and
\[
 \mathrm{e}^{\lambda a}f\,\mathrm{d}V
 \leq cf\,\mathrm{d}V=\omega_a^n=\omega_b^n
 \leq\mathrm{e}^{\lambda b}f\,\mathrm{d}V.
\]
Apply \cref{cor:bounded-exponential-comparison} first to $a,u_\lambda$ and then to $u_\lambda,b$. This gives $a\leq u_\lambda\leq b$, proving \eqref{eq:exp-barriers}. Multiplication by $\lambda$, together with $m\leq u\leq M$, gives the norm bound and the uniform limit.

If another pair $(v,c')$ solves the equation with the same $f$, the same unique functions $u_\lambda$ satisfy
$\lambda u_\lambda\to\log c'$ as well. Hence $c'=c$.
\end{proof}

The limiting characterization is due to Nguyen \cite[Corollary 0.2]{Nguyen16}. The displayed proof isolates its comparison mechanism. It proves uniqueness of the constant without comparing two candidate potentials for the equation with $\lambda=0$. The next chapter gives explicit examples and develops criteria for potential uniqueness. For classical potentials, uniqueness of the scalar also has a direct extremum proof:

% Not input: the statement is contained in prop:vanishing-exponent, which carries its label.


\section{Normalization on the Hopf threefold}
\label{sec:hopf-lp-normalization}

Let $X=(\mathbb C^3\setminus\{0\})/\langle z\mapsto\mathrm{e}^{-1}z\rangle$, with projection $\pi:X\to\mathbb P^2$ and the Hermitian form $\omega=\alpha+\beta$ of \cref{prop:hopf-mass}. Here $\alpha=\pi^*\omega_{\FS}$, $\int_{\mathbb P^2}\omega_{\FS}^2=1$, and $\nu=\omega^3/6$ is a probability measure. Let $E=\mathbb C/(\mathbb Z+2\pi\mathrm{i}\mathbb Z)$ act by $\tau_a[z]=[\mathrm{e}^a z]$, and let $A$ denote averaging with respect to its probability Haar measure $\mathrm{d}A_E$, as in \cref{prop:hopf-fiber-average}. The measure $\nu$ disintegrates as $\mathrm{d}A_E\,\omega_{\FS}^2$. For $u\in\PSH(X,\omega)$, write $\omega_u=\omega+\ddc u$.

Fix $p>1$ and a nonnegative nonzero $f\in L^p(X,\nu)$. By \cref{thm:kn-continuous,prop:vanishing-exponent}, there is a unique positive scalar $c_f$ for which
\[
 (\omega+\ddc u)^3=c_f f\,\omega^3
\]
has a continuous $\omega$-plurisubharmonic solution. Put $I_f=\int_Xf\,\mathrm{d}\nu$. The bounded mass characterization in \cref{prop:hopf-fiber-average} extends the smooth equality case of \cref{prop:hopf-normalization-defect} to these densities.

\begin{proposition}\label{prop:hopf-lp-normalization}
One has $c_f I_f\geq1$, and every continuous solution $u$ satisfies
\begin{equation}\label{eq:hopf-lp-average}
 \|u-Au\|_{L^2(X,\nu)}^2\leq4(c_f I_f-1).
\end{equation}
Moreover, $c_f I_f=1$ if and only if $f=\pi^*g$ almost everywhere with respect to $\nu$, for some nonnegative $g\in L^p(\mathbb P^2,\omega_{\FS}^2)$. In this case every continuous solution is the pullback of a continuous $\omega_{\FS}$-plurisubharmonic function.
\end{proposition}
\begin{proof}
Integrating the equation gives $\int_X\omega_u^3=6c_f I_f$. The lower mass bound and \eqref{eq:hopf-average-mass} give the first two assertions.

Suppose $c_f I_f=1$. By \cref{prop:hopf-fiber-average}, $u=\pi^*q$ for a bounded $\omega_{\FS}$-plurisubharmonic function $q$. Local holomorphic sections of $\pi$ show that $q$ is continuous. Thus $u\circ\tau_a=u$ for every $a\in E$. Naturality of bounded products and invariance of $\omega$ give
\[
 \tau_a^*(f\nu)=f\nu\qquad(a\in E).
\]
This is a measure identity; no common pointwise representative of $f$ for all $a$ is needed.

Fubini's theorem defines the fiber average $Af$ almost everywhere. It equals $\pi^*g$ for a measurable function $g$ on the base: local trivializations differ by translations of the Haar variable. Jensen's inequality gives
\[
 \int_{\mathbb P^2}|g|^p\omega_{\FS}^2
 \leq\int_X|f|^p\,\mathrm{d}\nu.
\]
For every continuous real test function $\chi$ on $X$, measure invariance gives $\int_X\chi f\,\mathrm{d}\nu=\int_X(\chi\circ\tau_a)f\,\mathrm{d}\nu$. Integrate in $a$ and use Fubini and invariance of Haar measure under $a\mapsto-a$. The result is $\int_X\chi f\,\mathrm{d}\nu=\int_X\chi Af\,\mathrm{d}\nu$. Hence $f=Af=\pi^*g$ almost everywhere, with $g\geq0$.

Conversely, suppose $f=\pi^*g$ almost everywhere. Disintegration gives $I_f=\int_{\mathbb P^2}g\omega_{\FS}^2$. Apply \cref{thm:kn-continuous} on the base. There are a continuous $\omega_{\FS}$-plurisubharmonic function $q$ and a scalar $k>0$ with
\[
 (\omega_{\FS}+\ddc q)^2=kg\,\omega_{\FS}^2.
\]
Bounded decreasing regularization and Stokes give $\int_{\mathbb P^2}(\omega_{\FS}+\ddc q)^2=1$, so $k=I_f^{-1}$.

Put $u=\pi^*q$. To compute its bounded product, take smooth strictly $\omega_{\FS}$-plurisubharmonic $q_j\downarrow q$ by \cref{thm:bounded-exact-regularization}. Their pullbacks decrease to $u$ and are strictly $\omega$-plurisubharmonic. For each $j$, the smooth cube identity gives
\[
 (\omega+\ddc\pi^*q_j)^3
 =3\pi^*(\omega_{\FS}+\ddc q_j)^2\wedge\beta.
\]
On a local trivialization, the right side is $6\mathrm{d}A_E$ times the base measure $(\omega_{\FS}+\ddc q_j)^2$: horizontal components of $\beta$ vanish against the base top-degree form, and its vertical restriction is $2\mathrm{d}A_E$. Bounded-product continuity on the base and on $X$, tested in these trivializations and patched by a partition of unity, therefore gives
\[
 (\omega+\ddc u)^3=\frac{f}{I_f}\,\omega^3.
\]
Uniqueness of the scalar implies $c_f=I_f^{-1}$. Every continuous solution then has mass $6$, so the same mass characterization makes it a base pullback. Its descended potential is continuous by local sections, as above.
\end{proof}

The argument allows densities that vanish and uses the bounded mass-minimizer theorem directly. It makes no assertion about a vertical-energy identity for nonsmooth potentials or about unbounded solutions.


\Needspace{17\baselineskip}
\section{Limits of this chapter}

Let $(X,\omega)$ be a connected compact Hermitian manifold and fix a smooth positive volume form $\mathrm{d}V$. For a nonnegative nonzero density $f\in L^p(X,\mathrm{d}V)$, $p>1$, denote by $u_\lambda$ the continuous solution of \eqref{eq:weak-exp}, with $\lambda>0$.

The continuous existence proofs use local weighted Chern--Levine--Nirenberg and volume--capacity estimates as standard inputs. The approximating equations control their own densities, so no bounded mass hypothesis on $X$ is needed.

The scalar limit above does not prove that the normalized potentials
$u_\lambda-\sup_Xu_\lambda$ have a unique limit for every density under consideration. \Cref{prop:symmetric-equation-clusters} uses the compactness theorem above to prove their uniform precompactness and preservation of data symmetries. Agreement of all cluster limits remains a further question.

For equations with prescribed singularities, we must also specify the reference type and the class of admissible solutions. The following chapters develop the relative envelopes and volume conditions used for this purpose, including the results of \cite{ALS26,XL26}.
